% ============================================
% Assistant Professor Cover Letter – MSU Denver
% ============================================

\documentclass[11pt,letterpaper]{letter}

\usepackage{geometry}
\usepackage{microtype}
\usepackage[hidelinks]{hyperref}

\geometry{left=25mm,right=25mm,top=25mm,bottom=25mm}

\signature{Decory Edwards}
\address{
Department of Economics \\
Johns Hopkins University \\
Baltimore, MD 21218 \\
\texttt{dedwar65@jh.edu}
}

\begin{document}

\begin{letter}{Search Committee \\
Department of Economics \\
Metropolitan State University of Denver \\
Denver, CO \\ USA}

\opening{Dear Members of the Search Committee,}

I am writing to apply for the Assistant Professor position in Economics at Metropolitan State University of Denver. I am a Ph.D. candidate in Economics at Johns Hopkins University, advised by Professors Christopher D.~Carroll, Nicholas W.~Papageorge, and Stelios Fourakis, and I expect to receive my degree in May 2026. My research fields are macroeconomics, wealth and income dynamics, and computational economics.

My research agenda focuses on understanding the determinants of wealth inequality and the mechanisms through which heterogeneity in returns to wealth shapes aggregate outcomes. In my job-market paper, I estimate the degree of heterogeneity in individual portfolio returns using micro-data from the Survey of Consumer Finances and simulate a structural model in which households face idiosyncratic return risk. I show that allowing for persistent heterogeneity in returns substantially improves the model’s ability to match the empirical distribution of wealth, offering a tractable way to connect micro-level return dispersion to macroeconomic inequality.

In a second project, I use the Health and Retirement Study to examine the relationship between interpersonal trust and portfolio performance. Building on the literature that finds a hump-shaped relationship between trust and income, I show that a similar relationship holds for individual returns to wealth measured in the HRS data, highlighting a behavioral channel through which social attitudes shape saving and investment outcomes. This work also provides new estimates of the persistent component of returns to net worth, which motivate the calibration of heterogeneity in my structural model.

MSU Denver is an especially attractive setting for my work and teaching. The university’s mission of access and student success, together with the distinctive strengths of an urban public institution, align with my interest in research and instruction that connect rigorous quantitative analysis to real economic outcomes. MSU Denver’s designation as a Hispanic-Serving Institution and its role in serving a diverse, largely transfer and first-generation student community are particularly compelling to me. :contentReference[oaicite:1]{index=1} I would be excited to contribute to a department that values high-quality teaching, supports student mobility, and prepares students to use economic reasoning and data in professional and civic life.

\textbf{Teaching.} I have experience teaching and supporting courses in \textit{Principles of Macroeconomics}, \textit{Intermediate Macroeconomic Theory}, and upper-level electives in macroeconomic policy. My teaching philosophy emphasizes clarity, economic intuition, and the use of data to connect theory to real-world outcomes. In a setting like MSU Denver, I would bring a structured, student-centered approach that meets students where they are, while building confidence with quantitative tools through frequent practice, transparent expectations, and applied examples. I am prepared to teach principles and intermediate macroeconomics, and I can contribute to economics courses that emphasize data analysis and computational methods.

Beyond academic fit, I am drawn to MSU Denver’s location and mission. I value working in a diverse, student-centered environment and I would welcome the opportunity to build a long-term academic career in Denver while contributing to MSU Denver’s goals in teaching, service, and scholarly engagement.

Thank you very much for your time and consideration. I would welcome the opportunity to discuss my application further.

\closing{Sincerely,}

\end{letter}
\end{document}
